% ============================================================
%  Event responses and the emitted product
%  How the network producer answers events on its carried state
%  and what it writes on the product wire, under the golden
%  configurations armed today.
%
%  Sources: filter/NetworkResponses, NetworkCarryDisposition,
%  NetworkClockAnchor, NetworkGrossLos, NetworkSolPrep,
%  NetworkPbHold, NetworkEmitStore; app/NetworkRunContext,
%  ApplicationNetwork; slip/NetworkSlipQc, NetworkRawObservables;
%  corrections/GnssCorrections; products/NetworkProductIO;
%  realtime/RealTimeEngineSsrBreaks, PbHoldLedger; config loaders
%  and posture rows; config/golden/golden_{field,sim}_refound.json.
%  Standalone-compilable subfile.
% ============================================================
\documentclass[main.tex]{subfiles}
\begin{document}
\ifSubfilesClassLoaded{\begin{center}\small\itshape Part of the GnssProcessor Mathematical Reference \textbullet\ \today\end{center}}{}

\section{Event responses and the emitted product}\label{sec:net-wire}

% Doc-local macros with distinctive names (subfiles collision note).
\providecommand{\wireWuHorizon}{T_{\mathrm{wu}}}  % wind-up continuity horizon
\providecommand{\wirePlHorizon}{T_{\mathrm{pl}}}  % leveling-carry gap horizon
\providecommand{\wireHoldGap}{T_{\mathrm{h}}}     % hold generation gap
\providecommand{\wireRdHorizon}{T_{\mathrm{rd}}}  % re-admission horizon
\providecommand{\wireWrap}{\operatorname{w}}
\providecommand{\wireMed}{\operatorname*{med}}

This section covers two duties of the network producer of
\dcref{sec:net}{the Network-mode document}: the response to events on
carried state (a slip, a tracking gap, a satellite's absence and return,
a restart of the wind-up model, a gross code fault on one line of sight),
and the product wire (which phase-bias value reaches the consumer, how a
change is declared, how the consumer answers). The carried local column,
the bounded drop operator, the held S-basis and the coupling time update
belong to \dcref{sec:net-vehicle}{the float-ambiguity-vehicle section}.
Knobs are quoted with their loader defaults; ``the goldens'' are
\texttt{config/golden/golden\_field\_refound.json} and
\texttt{golden\_sim\_refound.json}, which agree on every knob named here.

\subsection{The event doctrine}\label{sec:net-wire-doctrine}

A carried quantity enters the next epoch as a prior. In the information
form of \dcref{sec:net-carry}{the between-epoch carry document}, a
carried value $\hat x$ with variance $P$ folds as the pseudo-observation
\begin{equation}
x = \hat x + e, \qquad e \sim \mathcal{N}\big(0,\; P + q\,\dt\big),
\label{eq:wire-prior}
\end{equation}
with $q$ the process-noise rate of its kind. If the quantity moved by
$\delta$ after $\hat x$ was formed, the pseudo-observation is biased by
$\delta$ at weight $(P + q\,\dt)^{-1}$. For a time-linked quantity
$q \approx 0$ and $P$ is small, so the states sharing its rows
(ionosphere, clocks, troposphere) absorb the bias; a wired response must
change $(\hat x, P)$ before the fold. Standing invariant~7 of
the project states the rule: a mechanism that changes which
information flows into which states is an event at both edges. The
decoupling edge bounds what the decoupled states can accumulate; the
re-coupling edge re-validates or re-seeds what it reconnects. Each edge
writes an event row and a counter, so an inert response reads as zero.

\paragraph{Vocabulary.}
Card PRM phase~3 (\texttt{filter/NetworkResponses},
\texttt{NetworkParamRegistry}) names five events (\texttt{ParamEvent}:
slip, gap, handover, decouple, readmit) and eighteen policies
(\texttt{ResponsePolicy}). Those used here: \emph{Drop} returns the
variance to birth strength, clears the border regression and the block
row, and keeps the value (\dcref{sec:net-vehicle}{the
float-ambiguity-vehicle section}); \emph{EraseAndRebirth} erases the
entry for the next harvest to re-create; \emph{Park},
\emph{SeedAtVarPlusQGap} and \emph{Discard} keep an entry past its
column's departure, seed the returned column, or drop the entry
unfolded; \emph{Hold} freezes and re-applies an entry. Each response
site is one of 69 rows of \texttt{ResponseSite}, a (kind, event, policy,
code site) tuple with an arming predicate (\texttt{responseSiteArmed})
and a counter; a site that responds where it is not declared armed is a
run-end declaration error. Marks that only re-seed a linearization point
are excluded, among them the pending-reseed mark, which PASS~B's
measurement pre-pass consumes (a solve no-op on the float vehicle). The
acting ambiguity response is the free-epoch latch, which ends in a Drop
(card PRM3, section~2).

\paragraph{Armed set.}
The goldens arm the MW/LLI and IF triple-difference slip drains
(\texttt{network.cycle\_slip.mw\_detect}, \texttt{if\_phase\_td};
\dcref{sec:slip-network-qc}{the network slip-QC section}) with three
refinements: \texttt{mw\_evidence\_refloat} admits an MW flag on its own
strength (bias-free lane, mature arc, ratio to the threshold above
\texttt{mw\_evidence\_ratio\_min}); \texttt{common\_band\_scope} frees
only the bands the evidence implicates; \texttt{anchor\_slip\_handover}
bans a freed phase anchor's family for that epoch, so the election picks
a successor pinned at its own carried value. They also arm the responses
of \Cref{sec:net-wire-gap} to \Cref{sec:net-wire-consumer}. Each arming
knob is off in the loader.

\subsection{Gap re-entry, keyed per signal}\label{sec:net-wire-gap}

An arc that returns within $\wirePlHorizon$
(\texttt{network.phase\_leveling\_max\_gap\_s}, loader default
\SI{120}{s}, goldens \SI{900}{s}) has its carried entry re-injected at
the pre-gap value. No time-difference detector covers the return epoch:
MW needs an adjacent predecessor on the same pair, the IF triple
difference three consecutive epochs. The gap re-entry drains
(\texttt{slip/NetworkSlipQc}) run serially in the PASS-A slip band, after
the wind-up restart drain and before the inject, so an arc is freed in
the epoch it returns.

\paragraph{Row-level verdict.}
\texttt{networkComputeRawObservables} builds one raw row per (station,
satellite) on the two lowest distinct bands $s_1, s_2$ with phase and
matching code. A row absent for at least one processed epoch and present
now is a \texttt{GapResumption}, with verdict
\begin{equation}
\text{slipped} = \underbrace{\big(|\Delta\mwc| > g_{\mwc} \vee |\Delta\gfc| > g_{\gfc}\big)}_{\text{shifted}}
\vee \big(\kappa_{\mathrm{lli}} \wedge (\mathrm{LLI}_{s_1} \vee \mathrm{LLI}_{s_2})\big)
\vee \big(\kappa_{\mathrm{unm}} \wedge \neg\,\text{measurable}\big),
\label{eq:wire-gap-verdict}
\end{equation}
with $\Delta\mwc, \Delta\gfc$ the raw level changes across the break,
$g$ the station's empirical gates, $\kappa_{\mathrm{lli}}$ =
\texttt{gap\_reentry\_relevel\_lli} (default true) and
$\kappa_{\mathrm{unm}}$ = \texttt{gap\_reentry\_relevel\_unmeasurable}
(default false: missing evidence is not evidence of a slip). A slipped
row frees both bands.

\paragraph{The target is the carry.}
A returning arc has no live ambiguity column at drain time: it was pruned
with the gap, and PASS~B's pre-pass re-creates it after the drains. The
response synthesizes the inject's key (\texttt{StateKeys::makeAmbiguityKey})
and asks the carry whether anything needs protection: a leveling-carry
entry $(\rx,\sat,\sigl)$ with
\begin{equation}
0 \le t - t_{\mathrm{seen}} \le \wirePlHorizon ,
\label{eq:wire-protect}
\end{equation}
exactly what the inject would fold. A protectable arc enters the
free-epoch latch (\texttt{relevel.free\_ambs}, cause
\texttt{kCauseGap}); a live column also takes the pending-reseed mark.
\texttt{network.cycle\_slip.relevel\_free\_epoch} is the latch: it carries
the freed key past the pre-pass to the inject, where the arc folds at
birth strength after the bounded drop of \dcref{sec:net-vehicle}{the
float-ambiguity-vehicle section}. An arc without a protectable entry is
counted \texttt{no\_carry}.

\paragraph{Per-signal key.}
The raw row is rebuilt every epoch from whichever phases have a matching
code. A receiver that loses lock on its primary pair while another signal
tracks through (band~5, or a fallback signal) keeps its row, so the
row-level tier sees no gap and the returning phases would fold across the
break. \texttt{network.cycle\_slip.gap\_reentry\_per\_signal} (card GRS)
keys the re-entry on (station, satellite, signal). For a signal $\sigl$
last seen at epoch $k_\sigl$, the level test differences raw phases
(metres) against a partner $p$ present at $k$ and $k_\sigl$:
\begin{equation}
d_\sigl
= \big[\phaseo_\sigl - \phaseo_p\big](k) - \big[\phaseo_\sigl - \phaseo_p\big](k_\sigl)
\approx \lambda_\sigl\,\slipd N_\sigl - \lambda_p\,\slipd N_p - (\mu_\sigl - \mu_p)\,\delta I ,
\label{eq:wire-sig-dgf}
\end{equation}
up to noise, multipath and the wind-up change. A partner that tracked
through the gap is preferred: then $\slipd N_p = 0$ and $d_\sigl$
isolates the candidate's step plus the ionospheric change over the gap.
With $g_{\gfc}$ the station's GF gate, which does not grow with the gap,
\begin{equation}
\text{slipped}_\sigl = \big(|d_\sigl| > g_{\gfc}\big) \vee
\big(\kappa_{\mathrm{lli}} \wedge \mathrm{LLI}_\sigl\big) \vee
\big(\kappa'_{\mathrm{unm}} \wedge \neg\,\text{measurable}_\sigl\big),
\label{eq:wire-sig-verdict}
\end{equation}
with $\mathrm{LLI}_\sigl$ the returning signal's own flag and
$\kappa'_{\mathrm{unm}}$ = \texttt{gap\_reentry\_per\_signal\_unmeasurable}
(default true: the carry of a satellite that stayed present is certain
to fold). Only $\sigl$ is freed; the other bands tracked through. Both
tiers share one responder (\texttt{gapReentryRespondArc}); the
per-signal tier skips a satellite the row-level tier reported.

\paragraph{Tracker naming.}
Before card GRK the row-level tier named each band lexically: where the
canonical band-1 phase is L1W the tracker named L1C, the key missed the
carry, and the stale ambiguity folded.
\texttt{network.cycle\_slip.tracker\_modeled\_first} names each band by
the station's canonical phase latch (\texttt{TrackerLatchLookup}) when
that signal has raw phase and code this epoch, so the tracker key equals
the carried key; \texttt{gap\_tracker\_key\_mismatch} is then zero by
construction.

\subsection{Re-admission after a long absence}\label{sec:net-wire-readmit}

A satellite in eclipse, or below every mask, leaves the shared layout,
and its clock and code-bias columns depart; on return they re-realize
from code. \texttt{network.readmit\_seed\_horizon\_s} $=\wireRdHorizon$
(loader default 0, off; goldens \SI{3900}{s}) bridges the absence.

At departure the stash (\texttt{app/NetworkRunContext}) parks each D2
clock entry whose column left the layout: it keeps the entry with its
departure time $t_{\mathrm{park}}$, unfolded. Code-bias entries park too
under \texttt{network.readmit\_seed\_code\_bias} (default true).
Phase-bias entries do not: their frame is re-realized at the cold re-pin
on return, and a stiff seed would conflict with that pin. When the column
returns with $0 < t - t_{\mathrm{park}} \le \wireRdHorizon$, the fold
decision (\texttt{carryFoldDecide}) is \emph{Seeded},
\begin{equation}
P_{\mathrm{seed}} = P + q\,(t - t_{\mathrm{park}}),
\label{eq:wire-seed}
\end{equation}
with $q$ the registry rate of its kind. A clock under card Q3's temporal
anchor (\texttt{network.clock\_temporal\_anchor}, armed) folds through
the anchor with the staleness bound raised to
$\max(\texttt{clock\_anchor\_max\_gap\_s}, \wireRdHorizon)$
(\texttt{clockAnchorFoldVariance}), so it folds anchored over its true
gap instead of re-drawing from code. The anchor's guard is the
re-coupling check: a trip within its verify window after a seed is
recorded as the data rejecting the pre-gap realization
(\texttt{readmit\_seed\_moved}) and decouples the clock.

A parked entry is dropped past the horizon (\texttt{ClkParkExpire},
\texttt{CbParkExpire}), when its column returns without consuming it
(\texttt{ClkParkDiscard}), and at a clock-datum change of its
constellation that the handover S-transform did not re-express
(\texttt{ClkParkDiscardDatum}). The returning satellite is then a
different arc and registers cold. The clock-only variant
(\texttt{readmit\_seed\_code\_bias} false) was falsified on the
triple-frequency sim bed of card READMIT: the clock tails returned. The
code-bias seeds fix the split between clock and code bias at the return;
without them the cold code-bias realization relocates into the clock
(\texttt{docs/cards/network-vehicle/READMIT-EDGE.md}).

\subsection{Carrier-phase wind-up: accumulation and its restart}\label{sec:net-wire-windup}

The wind-up angle $\theta^{\mathrm{raw}} \in (-\pi,\pi]$ is an
$\operatorname{atan2}$ (\dcref{sec:obsphys-windup}{the wind-up section of
the observation-physics document}), applied as $\lambda_\sigl\theta/2\pi$
on every phase signal of the satellite. In \emph{fractional} mode
(\texttt{estimation.accumulate\_phase\_windup} false, the loader default)
$\theta = \theta^{\mathrm{raw}}$, and each crossing of $\pm\pi$ steps the
applied correction by one cycle on every signal inside a continuous arc.
The goldens accumulate: \texttt{GnssCorrections::computePhaseWindUp}
keeps $(\theta_{k-1}, t_{k-1})$ per satellite and forms
\begin{equation}
\theta_k =
\begin{cases}
\theta_{k-1} + \wireWrap\big(\theta^{\mathrm{raw}}_k - \theta_{k-1}\big), & 0 < t_k - t_{k-1} < \wireWuHorizon,\\[2pt]
\theta^{\mathrm{raw}}_k \quad \text{(restart)}, & \text{otherwise},
\end{cases}
\qquad \wireWrap(x) = x - 2\pi \operatorname{round}(x/2\pi),
\label{eq:wire-unwrap}
\end{equation}
with $\wireWuHorizon$ = \texttt{estimation.phase\_windup\_max\_gap\_s}
(loader default \SI{300}{s}; not set by the goldens). A repeated call
within one epoch ($t_k = t_{k-1}$) reuses the accumulated $\theta_k$ and
is not a restart. A restart is recorded in the ledger
\texttt{wuRestarts}. The state is keyed by
satellite; each station owns its correction object, so the key is
injective.

\paragraph{The restart is an integer, band-common step.}
By \eqref{eq:wire-unwrap}, $\theta_k - \theta^{\mathrm{raw}}_k \in
2\pi\mathbb{Z}$ at every epoch, so the continuation a restart replaces
equals $\theta^{\mathrm{raw}}_k + 2\pi n$ for an integer $n$. The restart
moves the modelled phase of every signal by $-n\lambda_\sigl$: the same
integer number of cycles on every signal, which the float ambiguities see
as a slip $\slipd N_\sigl = n$ for all $\sigl$. The raw observations did
not change, so no raw-observable detector can see it. In residual space
\begin{equation}
\frac{f_1 n\lambda_1 - f_2 n\lambda_2}{f_1 - f_2}
= n\,\frac{\sLight - \sLight}{f_1 - f_2} = 0,
\qquad
c_1 n\lambda_1 - c_2 n\lambda_2 = n\,\frac{\sLight}{f_1+f_2} = n\,\lamnl ,
\label{eq:wire-mw-blind}
\end{equation}
with $c_1, c_2$ the coefficients of \dcref{eq:if-coeffs}{the IF
coefficients}: the step is MW-blind and appears in the ionosphere-free
combination as $n\lamnl = n\lamif$. It also happens across a gap of at
least $\wireWuHorizon$, where no time difference has an adjacent
predecessor. The leveling carry survives gaps up to $\wirePlHorizon$, and
$\wireWuHorizon < \wirePlHorizon$ on the goldens; a return after a gap in
$[\wireWuHorizon, \wirePlHorizon]$ therefore steps the model while the
carried ambiguity re-injects at its pre-gap value (invariant~1).

\paragraph{The restart response.}
\texttt{network.cycle\_slip.windup\_restart\_refloat} (disarmed by the
loader without accumulation or phase leveling) drains each present
station's ledger in the PASS-A slip band. A record not stamped with the
current epoch is refused as stale; a gap beyond $\wirePlHorizon$ is
\texttt{beyond\_carry}, since the carry was pruned with it. Otherwise
\texttt{networkWindupRestartTargets} resolves the stepped arcs: the
satellite's live ambiguity columns at the station if any exist, else its
leveling-carry entries satisfying \eqref{eq:wire-protect}. Each target
enters the free-epoch latch (cause \texttt{kCauseWindup}; the
pending-reseed mark only on a live column), and the satellite's open
bridge pends are voided. With \texttt{cycle\_slip.refloat\_band\_scope}
at its default \texttt{all} every target band is freed; the
accumulator's verdict needs no raw corroboration. Card WUP fixed the
drain, which walked only the live ambiguity map: a restarted arc's column
was pruned with the gap, so no in-horizon restart reached an arc, the
inject folded the carry across the step, and the counter reported an
answer.

\paragraph{One mode for producer and consumer.}
The measured gain of accumulation is in the consumer's own rows; the wire
alone is neutral (card WUP). The wire declares the wind-up construction
(S-line \texttt{windup\_id} = \texttt{nominal-wu93}) and not the mode.
The in-process consumer inherits the producer's mode through the derived
configuration; a replay or stand-alone consumer holds the rule by
configuration.

\subsection{The gross line-of-sight screen}\label{sec:net-wire-gross}

Post-fit RAIM (\dcref{sec:stoch-net-raim}{the producer-side RAIM
section}) admits a row only if it carries an ambiguity coefficient above
0.5. A code row carries none, so a gross code error on one line of sight
never enters $\hat\sigma_0$, a cell maximum or the median polish.
\texttt{network.gross\_los\_screen} (\texttt{filter/NetworkGrossLos}) is
the acting tier for code rows; card RF-3 armed it with two sub-rules.

\paragraph{Statistic.}
Per station the seam forms $\bm{v} = \bm{y} - \bm{A}_s\hat{\bm{x}}_s -
\bm{A}_r\hat{\bm{x}}_r$ with $\hat{\bm{x}}_r$ from
\dcref{eq:net-backsub}{the back-substitution of the Network-mode
document}. A code row touches no SH, ambiguity or \texttt{PB:} column,
touches a \texttt{CB:} or \texttt{CLK:} column and names a satellite;
coupling rows are skipped. The statistic
\begin{equation}
r_j = |v_j|\,\sqrt{w_j}\; f_j
\label{eq:wire-gl-stat}
\end{equation}
uses the row's own weight $w_j$ and undoes the deweight factor $f_j$
the assembly applied (1 if none), so both bars read the data at nominal
weight. The row's own $\sigma$ is used because the pooled $\hat\sigma_0$
inflates while a station degrades, exactly when the screen must stay
sensitive.

\paragraph{Declaration and release.}
For one (station, satellite) with $n$ code rows, $n_{\mathrm{over}}$ of
them above $k$ = \texttt{network.gross\_los\_k}, the epoch's evidence
clears the bar when
\begin{equation}
n_{\mathrm{over}} \ge m_{\mathrm{eff}}, \qquad
m_{\mathrm{eff}} =
\begin{cases}
\max\big(1, \min(m, n)\big) & \texttt{available},\\
\max(1, m) & \texttt{fixed},
\end{cases}
\label{eq:wire-gl-bar}
\end{equation}
with $m$ = \texttt{network.gross\_los\_min\_rows} (loader default 2) and
\texttt{network.gross\_los\_min\_rows\_rule} (loader default
\texttt{fixed}). A gross line of sight is wrong on every code band at
once; one heavy-tailed row is one band's multipath. Under \texttt{fixed}
a line of sight with one code row can never be declared;
\texttt{available} reaches it and still demands two rows where two
exist. Declaring takes \texttt{gross\_los\_epochs} consecutive epochs of
evidence (absence resets the streak) and passes two counted guards: the
maturity horizon (\texttt{gross\_los\_warmup\_epochs}) and the
whole-station stand-down (\texttt{gross\_los\_max\_station\_frac}: a
station-wide fault is not a line-of-sight fault). Release takes every
code row under \texttt{gross\_los\_clear\_sigma} for
\texttt{gross\_los\_clear\_epochs} epochs and is never guarded. A
declaration unseen beyond \texttt{gross\_los\_max\_gap\_s} expires with
its own event: the next sighting is a new arc.

\paragraph{Response.}
From the next assembly the declared code rows fold at
$\sigma\times$\texttt{gross\_los\_deweight}
(\texttt{networkGrossLosPublishScales}); the row keeps its sign and still
anchors the float if the phase side disappears. At the declaration the
line of sight's leveling-carry entries (\texttt{AmbGrossErase}, with a
pending-reseed mark) and slant-ionosphere entries
(\texttt{IonoGrossErase}) are erased. Both sites are declared
EraseAndRebirth on a decouple event: this response erases where the slip
responses drop. Phase rows keep their weight; they have their own tiers.
The bounded variant \texttt{gross\_los\_bound} (freeze instead of erase)
was falsified and stays off
(\texttt{docs/refound/RF-3\_code\_row\_qc\_2026-09-22.md}).

\paragraph{Two halves and the emission exclusion.}
The erase must follow the carry harvest (\texttt{networkLeveling}), or
the harvest re-stashes the erased entries; the decision needs only the
solved epoch. With \texttt{network.gross\_los\_emit\_exclude} the seam is
split (\texttt{GrossLosPhase}): the Decide half runs immediately before
\texttt{networkSolPrep}, the Act half after \texttt{networkPostFitRaim}.
Unsplit, a declaration reaches the wire one epoch late. The emitted clock
and code bias are joint-solve states, and removing a station's
contribution would take a second solve, so the exclusion acts on the
records. For each satellite $\sat$ screened at $n_\sat$ stations, $e_\sat$
of them with a standing declaration on it,
\begin{equation}
n_\sat - e_\sat < N_{\min} \;\Rightarrow\; \text{withhold } \sat;
\qquad \text{otherwise emit } \sat \text{ with a discontinuity mark},
\label{eq:wire-gl-exclude}
\end{equation}
with $N_{\min}$ = \texttt{network.gross\_los\_sat\_min\_stations}
(loader default 3). The withhold is per satellite and the epoch still
emits. The mark joins the retro queue, whose drain this knob arms, and
lands with full channel scope (clock, code bias, phase bias) in the
first emission containing the satellite, normally the declaring epoch.

\subsection{The emitted phase bias}\label{sec:net-wire-pb}

For satellite $\sat$, signal $\sigl$ and emission epoch $t$, let
$c^\sat(t)$ be the emitted clock ($C_0$, metres) and $p^\sat_\sigl(t)$
the estimator's phase bias in the emit frame (the \texttt{bias\_m}
\texttt{networkPbHold} receives). The wire carries $C^\sat = c^\sat$ and
$P^\sat_\sigl$; the consumer applies $-C$ on code and $-C + P$ on phase
(\dcref{eq:ssr-phase-bias}{the SSR phase-bias equation}).

\paragraph{The generation hold.}
Under \texttt{network.emit\_pb\_hold} (card PBH) each key $(\sat,\sigl)$
carries a generation $g \ge 1$, started at $t_0$ when the key is first
emitted or was last emitted more than $\wireHoldGap$ ago
(\texttt{network.emit\_pb\_hold\_max\_gap\_s}, loader default
\SI{300}{s}; not set by the goldens). Within a generation
\begin{equation}
P^\sat_\sigl(t) = p^\sat_\sigl(t_0).
\label{eq:wire-hold-plain}
\end{equation}
The generation is declared by a \texttt{gen=}$g$ token on every P record,
by \texttt{SsrPhaseBiasEntry::hold\_generation} in the in-process store,
and by \texttt{pbhold=} on the S line (1 plain, 2 phase-clock mode). The
solve and the C and B records are byte-identical with the knob on. The
premise, graded against truth in card PBH: the wire's phase-bias motion
within an arc is error that adds to the clock error. Like the
\texttt{iode=} cutover, the contract states that the value has not
changed and names the epoch at which that statement restarts; handover
steps of the estimate stay off the wire.

\paragraph{Ordering contract.}
\texttt{networkPbHold} is the last writer of \texttt{bias\_m} in
\texttt{networkEmitStore}: after the carry-invariant emission re-base and
the phase-bias emission filter (both gated on
\texttt{network.carry\_invariant}, inert on the goldens) and after every
discontinuity-marking site; before the emitter's value-jump test, which
then screens the held value; before the file write, which reads the same
solution as the in-process store.

\paragraph{Both edges.}
Decoupling: each held emission records $|p^\sat_\sigl(t) -
P^\sat_\sigl(t)|$ into a run-end distribution. Re-coupling: each
generation start is a counted \texttt{PB\_HOLD} \texttt{gen\_start} event
(cause \texttt{first} or \texttt{gap}). A mark naming a satellite's clock
and phase bias together states that both moved; the plain hold drops the
phase-bias half and records the de-commoned clock step it left out (the
phase-clock mode below transfers it).

\paragraph{The phase-clock mode.}
Under \texttt{network.emit\_pb\_hold\_phase\_clock} (card LAH, armed by
card EXL) the common clock motion of constellation $k$ accumulates over
consecutive emissions $t^- < t$,
\begin{equation}
M_k(t) = M_k(t^-) + \wireMed_{j \in \mathcal{S}_k(t)} \big[c^j(t) - c^j(t^-)\big],
\label{eq:wire-cm}
\end{equation}
over the satellites $\mathcal{S}_k(t)$ of $k$ with a finite clock at both
emissions (an even count takes the mean of the two middle values;
\texttt{pbHoldCommonClockStep}). With $\tilde c^\sat(t) = c^\sat(t) -
M_{k(\sat)}(t)$ the held emission is
\begin{empheq}[box=\fbox]{equation}
P^\sat_\sigl(t) = p^\sat_\sigl(t_0) + \big[\tilde c^\sat(t) - \tilde c^\sat(t_0)\big].
\label{eq:wire-hold-pc}
\end{empheq}
An emission without the satellite's clock reuses the last transfer; a
generation started without one is held plainly. Subtracting from
$C^\sat = c^\sat$ gives the phase-side identity
\begin{equation}
C^\sat(t) - P^\sat_\sigl(t)
= c^\sat(t) - p^\sat_\sigl(t_0) - c^\sat(t) + M_k(t) + c^\sat(t_0) - M_k(t_0)
= \big[c^\sat - p^\sat_\sigl\big](t_0) + M_k(t) - M_k(t_0).
\label{eq:wire-identity}
\end{equation}
The constant goes into the consumer's float ambiguity. The common mode is
part of every $c^j$ of the constellation, so it enters code and phase
rows alike and the receiver clock and inter-system bias absorb it.
Satellite-specific clock motion reaches the code rows only. In the family
$(C - P)(t) = (c - p)(t_0) + \Phi(t)$ the phase-clock mode is
$\Phi = M_k(t) - M_k(t_0)$; the plain hold is $\Phi = c(t) - c(t_0)$,
every clock move on the phase side; no hold passes the estimator's own
phase-visible change, wander included.

\paragraph{Why the mode exists.}
Under the exact coupling time update armed on the goldens
(\texttt{network.carry\_coupling\_exact\_update}, card CTU;
\dcref{sec:net-vehicle}{the float-ambiguity-vehicle section}) the
satellite's level converges: the estimator moves $c^\sat$ and every
$p^\sat_\sigl$ by the same amount, invisible to phase. The plain hold
freezes $P$ while $C$ follows, turning the convergence into a phase-side
drift the consumer's float ambiguities cannot follow. Under the
phase-clock mode the move shifts $\tilde c^\sat$, hence $P$, by the same
amount (unless common to the constellation): $C - P$ stays and the level
reaches the code side only.

\paragraph{Limitation.}
The mode is interim. Within a generation the phase-visible clock motion
the consumer sees is the a priori clock product's alone, since the
emitted clock refines the a priori clock (\dcref{sec:net-closedloop}{the
closed-loop section}) and the hold passes none of the network's
satellite-specific motion to the phase side. That is valid with a precise
a priori clock (\texttt{network.apriori\_source} \texttt{sp3}, the
default). The loader warns under \texttt{brdc\_rts}, refuses the mode
under \texttt{sp3\_brdc\_clock}, and refuses it with
\texttt{network.clock\_level\_carry} (both move the held bias with the
level) and with \texttt{network.fast\_clock.enable} (sub-epoch clock
records would saw-tooth the phase-visible clock). The smoothed members
(\texttt{network.emit\_pb\_hold\_smoothed}) are built and not armed.

\begin{remark}
A datum move between two generation starts, such as a phase-anchor
handover with its S-transform, moves the estimator's phase biases and not
the held values (counted, measure-only:
\texttt{prm\_pbhold\_unmoved\_at\_shift}). A float consumer absorbs each
held constant in its own ambiguity. An integer-resolving consumer needs
every phase bias of a constellation in one S-basis
(\dcref{sec:net-phasebias}{the explicit phase-bias section};
\cite{odijk2016estimability}), which a held wire does not guarantee after
a datum move. The goldens run the consumer float
(\texttt{closed\_loop.ar\_mode} 0).
\end{remark}

\subsection{The consumer's answer}\label{sec:net-wire-consumer}

The in-process consumer (undifferenced on the goldens) reacts in
\texttt{realtime/RealTimeEngine\allowbreak SsrBreaks}, in a fixed order.
First,
\texttt{reactToGenerationChange} reads the product's solution-generation
id, which is distinct from the hold generation; a change (a producer
restart, a global re-realization) resets every network-corrected
ambiguity and re-inflates the receiver code-bias states. It runs before
the per-signal reactions because a restart re-seeds the discontinuity
counters.

\paragraph{Hold generation.}
\texttt{resetOnPhaseBiasDisc} dispatches \texttt{reactToPbHold\_} before
any other phase-bias reaction, keyed on the data: a key whose store entry
has a nonzero \texttt{hold\_generation} is claimed. Its discontinuity
counter is consumed, so the per-key memory advances, and ignored, since
under a hold a counter change carries no value step.
\texttt{PbHoldLedger::react} decides: \emph{FirstSight} adopts the
generation, \emph{Held} keeps it, \emph{Refloat} answers a change of a
previously seen generation with \texttt{StateManager::resetAmbiguity}
at the configured ambiguity $\sigma$, the call a discontinuity reset
uses. First sightings, held cells and refloats sum to the visited cells.

\paragraph{Presence edges.}
Under the bias staleness gate (\texttt{processing.enable\_ssr\_bias\_max\_age},
armed by the closed-loop harness on its derived consumer) the consumer
tracks
\begin{equation}
a_\sigl(t) = \big[\text{bias present} \wedge \text{age} \le
\texttt{ssr\_bias\_max\_age\_s} \cdot \texttt{ssr\_bias\_stale\_hold\_factor}\big]
\label{eq:wire-presence}
\end{equation}
and answers an edge $a_\sigl(t) \ne a_\sigl(t^-)$, the A5 freshness flip,
with an ambiguity reset. A held wire has two such edges with no
generation change: a held bias appearing on a phase row already in use
without it, and one leaving such a row; either steps the applied
correction by the whole bias. Before card PBF a claimed key skipped every
downstream reaction, A5 included, and its A5 memory was not written while
held, so both edges went unanswered (a first sighting lacks prior state
only on a new row). The update's post-fit peel then attributed the
residual to other satellites, reset their ambiguities and could exclude
the faulty satellite's code rows; card PBF records a field epoch that
moved the position by about \SI{490}{m}. With
\texttt{processing.ssr\_pb\_hold\_presence\_flip}, derived from
\texttt{network.closed\_loop.pb\_hold\_presence\_flip}, a claimed key
skips only the counter reaction, which the hold owns, and its A5 edge
runs as on an unheld key: on the goldens a reset, since the step
translation (\texttt{closed\_loop.pb\_step\_translate}) and the level
check (\texttt{closed\_loop.amb\_level\_check}, harmful under the hold,
card PBH) are off.

\subsection{Standing}\label{sec:net-wire-standing}

A card named by its identifier, such as PBH, is the file
\texttt{docs/refound/}\emph{identifier}\texttt{\_*.md}. Armed on both
goldens through the promotion ledger: gap re-entry per row
(\texttt{docs/cards/slip/CARD-G-GAP-REENTRY-RELEVEL.md}), per signal
(GRS) and with the tracker's band naming (GRK); the re-admission horizon
(\texttt{docs/cards/network-vehicle/READMIT-EDGE.md}); accumulated
wind-up with its restart response (WUP); the gross line-of-sight screen
(RF-3); the generation hold (PBH) in its phase-clock mode (LAH, armed by
EXL); the consumer's presence edges (PBF). Owed: the hold in the
warm-resume snapshot and the cause of the phase-bias wander (PBH); the
held emission under a phase-bias shift (PRM3); LAH's regressions at the
C1W stations and its need for a precise a priori clock (pursued by SCL
and NBP); RF-3's dedicated wire writer; GRS's failed field
triple-frequency bar; the peel's masking of a bad phase row (PBF). The
register is \texttt{docs/status/OPEN\_ITEMS\_A\_current\_program.md}.

\ifSubfilesClassLoaded{\printbibliography}{}
\end{document}
